\documentclass[12pt, a4paper]{article}

\usepackage{fontspec}
\usepackage{polyglossia}
\setmainlanguage{french}
\usepackage{amsmath, amssymb}
\usepackage{geometry}
\usepackage{booktabs}
\usepackage{listings}
\usepackage{xcolor}
\usepackage{hyperref}
\usepackage{parskip}
\usepackage{algorithm}
\usepackage{algpseudocode}
\usepackage{graphicx}

\geometry{margin=2.5cm}

\definecolor{codebg}{RGB}{245,245,245}
\definecolor{codegreen}{RGB}{0,128,0}
\definecolor{codegray}{RGB}{128,128,128}
\definecolor{codeblue}{RGB}{0,0,200}

\lstset{
  language=Python,
  backgroundcolor=\color{codebg},
  basicstyle=\ttfamily\small,
  keywordstyle=\color{codeblue}\bfseries,
  commentstyle=\color{codegreen},
  stringstyle=\color{orange},
  numberstyle=\tiny\color{codegray},
  numbers=left,
  stepnumber=1,
  frame=single,
  framerule=0.4pt,
  breaklines=true,
  showstringspaces=false,
  tabsize=4,
}

\title{
  \textbf{Integer Linear Programming (ILP)}\\
  \large Documentation algorithmique --- Travelling Salesman Problem
}
\author{Recherche Opérationnelle --- CESI}
\date{}

\begin{document}

\maketitle
\tableofcontents
\newpage

% =============================================================================
\section{Définition}

L'\textbf{Integer Linear Programming (ILP)} est une méthode de résolution \textbf{exacte}
du Travelling Salesman Problem (TSP). Le problème est modélisé comme un programme
mathématique : on définit des \textbf{variables binaires} représentant les arcs empruntés,
une \textbf{fonction objectif} à minimiser (distance totale), et des \textbf{contraintes}
garantissant un tour hamiltonien valide.

Un solveur (OR-Tools, CPLEX, Gurobi) explore l'espace des solutions via du
\textbf{branch-and-cut} : il divise récursivement l'espace de recherche tout en coupant
les branches sous-optimales grâce à des inégalités linéaires.

% =============================================================================
\section{Méthode commune aux deux formulations}

\subsection{Variables de décision}

\begin{equation}
  x_{ij} \in \{0, 1\}, \quad \forall\, i \neq j
\end{equation}

$x_{ij} = 1$ si l'arc de la ville $i$ vers la ville $j$ est emprunté, $0$ sinon.

\subsection{Fonction objectif}

\begin{equation}
  \min \sum_{i=1}^{n} \sum_{j \neq i} c_{ij} \cdot x_{ij}
\end{equation}

où $c_{ij} = \sqrt{(x_i - x_j)^2 + (y_i - y_j)^2}$ est la distance euclidienne.

\subsection{Contraintes de degré}

\begin{align}
  \sum_{j \neq i} x_{ij} &= 1 \quad \forall\, i \qquad \text{(quitter } i \text{ une fois)} \\
  \sum_{i \neq j} x_{ij} &= 1 \quad \forall\, j \qquad \text{(arriver en } j \text{ une fois)}
\end{align}

Ces contraintes seules ne suffisent pas : la solution peut se décomposer en plusieurs
cycles disjoints (sous-tours). Les formulations MTZ et DFJ traitent ce problème
différemment.

% =============================================================================
\section{Formulation MTZ (Miller--Tucker--Zemlin, 1960)}

\subsection{Principe}

MTZ introduit des variables auxiliaires $u_i$ représentant la
\textbf{position ordinale} de la ville $i$ dans le tour.
Elles forcent un ordre strict entre les villes visitées successivement,
ce qui rend tout cycle strict sur un sous-ensemble impossible.

\subsection{Variables supplémentaires}

\begin{equation}
  u_i \in \{0, 1, \ldots, n-1\}, \quad \forall\, i
\end{equation}

\subsection{Contraintes MTZ}

La ville de départ est fixée en position $0$ :

\begin{equation}
  u_0 = 0
\end{equation}

Pour tout arc emprunté $x_{ij} = 1$ avec $i \geq 1$ et $j \geq 1$, la ville $j$
doit être visitée strictement après $i$ :

\begin{equation}
  u_i - u_j + n \cdot x_{ij} \leq n - 1 \quad \forall\, i \neq j,\; i,j \geq 1
\end{equation}

\subsection{Preuve d'élimination des sous-tours}

Supposons un sous-tour $S = \{v_1, v_2, \ldots, v_k\} \subsetneq V$ avec $k < n$
et $v_1 \notin \{0\}$. En appliquant la contrainte MTZ à chaque arc du sous-tour :

\begin{align*}
  u_{v_1} - u_{v_2} &\leq -1 \\
  u_{v_2} - u_{v_3} &\leq -1 \\
  &\;\vdots \\
  u_{v_k} - u_{v_1} &\leq -1
\end{align*}

En sommant : $0 \leq -k$, ce qui est impossible pour $k \geq 1$.
Le sous-tour est donc interdit.

\subsection{Complexité}

MTZ ajoute $n$ variables entières et $O(n^2)$ contraintes linéaires.
Le modèle est compact mais la \textbf{relaxation LP est faible} : les variables $u_i$
peuvent prendre des valeurs fractionnaires sans signification physique,
obligeant le solveur à brancher massivement.

\begin{center}
\begin{tabular}{ll}
  \toprule
  Variables totales    & $O(n^2)$ binaires $+$ $n$ entières \\
  Contraintes totales  & $O(n^2)$ \\
  Limite pratique      & $n \leq 50$ \\
  \bottomrule
\end{tabular}
\end{center}

% =============================================================================
\section{Formulation DFJ avec Lazy Constraints (Dantzig--Fulkerson--Johnson, 1954)}

\subsection{Principe général}

Contrairement à MTZ qui introduit des variables d'ordre auxiliaires, DFJ
interdit les sous-tours \textbf{directement sur les variables $x_{ij}$}, sans
variable supplémentaire.

L'idée centrale : dans tout sous-ensemble $S$ de $k$ villes, si les arcs
internes à $S$ forment un cycle fermé, il y en a exactement $k$.
On interdit donc qu'il y en ait autant que $|S|$ --- au plus $|S|-1$ arcs
peuvent rester à l'intérieur de $S$.

\subsection{Contraintes DFJ}

Pour tout sous-ensemble strict non vide $S \subsetneq V$ avec $2 \leq |S| \leq n-1$ :

\begin{equation}
  \sum_{\substack{i \in S,\, j \in S \\ i \neq j}} x_{ij} \leq |S| - 1
  \label{eq:dfj}
\end{equation}

\textbf{Preuve que \eqref{eq:dfj} interdit les sous-tours.}
Supposons un sous-tour sur $S = \{v_1, \ldots, v_k\}$, c'est-à-dire
$x_{v_1 v_2} = \cdots = x_{v_k v_1} = 1$.
Le membre gauche vaut alors $k = |S|$, ce qui viole $|S| - 1$.

\textbf{Preuve que \eqref{eq:dfj} n'interdit pas les tours valides.}
Dans un tour hamiltonien, tout sous-ensemble $S$ de $k$ villes contient
exactement $k-1$ arcs internes. La contrainte est satisfaite avec égalité.

\subsection{Le problème de l'explosion combinatoire}

Il existe $2^n - 2$ sous-ensembles $S$ non triviaux, donc autant de contraintes
potentielles :

\begin{center}
\begin{tabular}{cc}
  \toprule
  $n$ & Contraintes DFJ potentielles \\
  \midrule
  10 & $2^{10} - 2 \approx 1\,000$ \\
  20 & $2^{20} - 2 \approx 1\,000\,000$ \\
  50 & $2^{50} - 2 \approx 10^{15}$ \\
  \bottomrule
\end{tabular}
\end{center}

Ajouter toutes ces contraintes dès le départ est hors de portée dès $n \geq 25$.

\subsection{Stratégie des lazy constraints}

On exploite le fait que la très grande majorité de ces contraintes ne sont
\emph{jamais violées}. On les génère donc \textbf{à la demande} :

\begin{algorithm}[H]
\caption{ILP DFJ avec lazy constraints (version itérative)}
\begin{algorithmic}[1]
\State $\mathcal{C} \leftarrow \emptyset$ \Comment{coupes DFJ actives}
\Repeat
  \State \textbf{Construire} le modèle ILP avec les contraintes de degré et $\mathcal{C}$
  \State \textbf{Résoudre} : obtenir $x^*$ optimal pour ce modèle réduit
  \State \textbf{Détecter} les sous-tours $\mathcal{S} = \{S_1, \ldots, S_k\}$ dans $x^*$
  \If{$|\mathcal{S}| = 1$}
    \State \Return $x^*$ \Comment{tour unique $\Rightarrow$ hamiltonien valide et optimal}
  \EndIf
  \For{chaque $S_\ell \in \mathcal{S}$ tel que $|S_\ell| < n$}
    \State $\mathcal{C} \leftarrow \mathcal{C} \cup \left\{\displaystyle\sum_{i \in S_\ell,\, j \in S_\ell,\, i \neq j} x_{ij} \leq |S_\ell| - 1\right\}$
  \EndFor
\Until{solution hamiltonienne trouvée}
\end{algorithmic}
\end{algorithm}

\textbf{Pourquoi la solution finale est optimale.}
À chaque itération, on résout le modèle courant à \textbf{l'optimalité exacte}.
Quand aucun sous-tour n'est détecté, toutes les contraintes DFJ nécessaires sont
satisfaites : $x^*$ est donc optimal pour le problème complet.

\subsection{Détection des sous-tours --- problème de séparation}

Étant donné $x^*$ entier, on construit le graphe $G^* = (V, \{(i,j) : x^*_{ij} = 1\})$.
Chaque composante connexe est un sous-tour, détectable par \textbf{DFS en $O(n)$} :

\begin{enumerate}
  \item Partir du sommet $0$, suivre les arcs $x^*_{ij} = 1$ jusqu'à revenir à $0$.
  \item Si tous les sommets sont visités $\Rightarrow$ tour hamiltonien.
  \item Sinon, chaque groupe non visité forme un sous-tour indépendant.
\end{enumerate}

Pour la relaxation LP (variables fractionnaires), un algorithme de
\textbf{coupe minimale} (min-cut) est nécessaire --- c'est le séparateur de Concorde.

\subsection{Convergence et nombre de coupes en pratique}

En pratique, seules \textbf{quelques dizaines à quelques centaines} de coupes
sont nécessaires, même pour des milliers de villes :

\begin{center}
\begin{tabular}{ccc}
  \toprule
  $n$ & Itérations typiques & Coupes ajoutées \\
  \midrule
  20  & 2--5  & 3--10  \\
  50  & 3--8  & 5--20  \\
  100 & 5--15 & 10--50 \\
  \bottomrule
\end{tabular}
\end{center}

\subsection{Comparaison MTZ vs DFJ}

La différence clé est la \textbf{qualité de la relaxation LP} :

\begin{itemize}
  \item \textbf{MTZ} : borne LP loin de l'optimum entier.
    Le solveur doit \textbf{brancher massivement} pour combler l'écart.
  \item \textbf{DFJ} : borne LP très proche de l'optimum.
    Concorde prouve souvent l'optimalité \textbf{sans brancher}, rien qu'avec les coupes.
\end{itemize}

\begin{table}[h]
  \centering
  \begin{tabular}{lcccc}
    \toprule
    \textbf{Formulation} & \textbf{Variables} & \textbf{Contraintes en mémoire} & \textbf{Qualité LP} & \textbf{Limite} \\
    \midrule
    MTZ        & $O(n^2)$ bin. $+$ $n$ ent. & $O(n^2)$ dès le départ  & Faible     & $n \leq 50$       \\
    DFJ (lazy) & $O(n^2)$ binaires          & Quelques centaines      & Excellente & $n \leq 100\,000$ \\
    \bottomrule
  \end{tabular}
  \caption{Comparaison détaillée MTZ vs DFJ}
\end{table}

% =============================================================================
\section{Branch-and-Cut}

Le branch-and-cut est le moteur interne des solveurs ILP modernes. Il combine
trois mécanismes :

\begin{enumerate}
  \item \textbf{Relaxation LP} --- On relâche $x_{ij} \in \{0,1\}$ en $x_{ij} \in [0,1]$.
        La solution LP fournit une \textbf{borne inférieure} en temps polynomial.

  \item \textbf{Cut} --- On cherche des inégalités valides (coupes) violées
        par la solution LP fractionnaire et on les ajoute au modèle.
        Coupes TSP principales : subtour elimination (DFJ), comb inequalities,
        blossom inequalities, Gomory cuts.

  \item \textbf{Branch} --- Si après coupes la solution reste fractionnaire
        ($x_{ij} = 0.7$), on crée deux sous-problèmes : $x_{ij} = 0$ et $x_{ij} = 1$.
        Les branches dont la borne LP dépasse la meilleure solution entière connue
        sont \textbf{élaguées} (pruning).
\end{enumerate}

% =============================================================================
\section{Explication intuitive}

On cherche à cocher exactement $n$ cases dans une grille $n \times n$
(une par ligne, une par colonne) de façon à former un \textbf{seul cycle} passant
par toutes les villes avec le coût minimum.

\medskip
\textbf{Exemple --- sous-tour sur 6 villes :}

\begin{center}
\begin{tabular}{cc}
  \toprule
  Solution invalide (2 sous-tours) & Solution valide (tour unique) \\
  \midrule
  $A \to B \to C \to A$ \quad $D \to E \to F \to D$ &
  $A \to B \to C \to D \to E \to F \to A$ \\
  \bottomrule
\end{tabular}
\end{center}

Les contraintes de degré acceptent les deux cas. MTZ les élimine via un ordre
global sur les villes ; DFJ les élimine en interdisant explicitement chaque cycle partiel.

% =============================================================================
\section{Cas d'applications connus}

\begin{table}[h]
  \centering
  \begin{tabular}{ll}
    \toprule
    \textbf{Référence} & \textbf{Contribution} \\
    \midrule
    Dantzig, Fulkerson, Johnson (1954)           & Premier algorithme exact TSP par coupes --- fondateur de DFJ \\
    Miller, Tucker, Zemlin (1960)                & Formulation compacte MTZ --- $O(n^2)$ contraintes \\
    Applegate, Bixby, Chvátal, Cook (Concorde)   & Solveur TSP spécialisé DFJ ; résolu jusqu'à 85\,900 villes \\
    OR-Tools CP-SAT (Google)                     & Moteur open source utilisé dans ce benchmark \\
    \bottomrule
  \end{tabular}
  \caption{Références principales}
\end{table}

% =============================================================================
\section{Exemple de code}

\subsection{Formulation MTZ}

\begin{lstlisting}[caption={ILP MTZ --- OR-Tools CP-SAT}]
from ortools.sat.python import cp_model
import math

def solve_tsp_mtz(cities, time_limit=120.0):
    n = len(cities)
    def dist(i, j):
        return int(math.hypot(cities[i][0]-cities[j][0],
                              cities[i][1]-cities[j][1]) * 1000)

    model = cp_model.CpModel()
    x = [[model.new_bool_var(f"x_{i}_{j}") for j in range(n)] for i in range(n)]
    u = [model.new_int_var(0, n-1, f"u_{i}") for i in range(n)]

    for i in range(n):
        model.add(x[i][i] == 0)
        model.add(sum(x[i][j] for j in range(n) if j != i) == 1)
        model.add(sum(x[j][i] for j in range(n) if j != i) == 1)

    model.add(u[0] == 0)
    for i in range(1, n):
        for j in range(1, n):
            if i != j:
                model.add(u[i] - u[j] + n * x[i][j] <= n - 1)

    model.minimize(
        sum(dist(i,j)*x[i][j] for i in range(n) for j in range(n) if i != j)
    )
    solver = cp_model.CpSolver()
    solver.parameters.max_time_in_seconds = time_limit
    solver.solve(model)
\end{lstlisting}

\subsection{Formulation DFJ avec lazy constraints}

\begin{lstlisting}[caption={ILP DFJ lazy --- CP-SAT iteratif}]
def find_subtours(n, x_val):
    visited, subtours = [False]*n, []
    for start in range(n):
        if visited[start]: continue
        tour, cur = [], start
        while not visited[cur]:
            visited[cur] = True
            tour.append(cur)
            cur = next(j for j in range(n) if x_val[cur][j] > 0.5)
        subtours.append(tour)
    return subtours

def solve_tsp_dfj(cities, time_limit=120.0):
    n = len(cities)
    def dist(i, j):
        return int(math.hypot(cities[i][0]-cities[j][0],
                              cities[i][1]-cities[j][1]) * 1000)

    forbidden = []
    while True:
        model = cp_model.CpModel()
        x = [[model.new_bool_var(f"x_{i}_{j}") for j in range(n)]
             for i in range(n)]
        for i in range(n):
            model.add(x[i][i] == 0)
            model.add(sum(x[i][j] for j in range(n) if j != i) == 1)
            model.add(sum(x[j][i] for j in range(n) if j != i) == 1)
        for S in forbidden:
            S = list(S)
            model.add(sum(x[i][j] for i in S for j in S if i != j) <= len(S)-1)
        model.minimize(sum(dist(i,j)*x[i][j]
                          for i in range(n) for j in range(n) if i != j))
        solver = cp_model.CpSolver()
        solver.parameters.max_time_in_seconds = time_limit
        solver.solve(model)
        x_val = [[solver.value(x[i][j]) for j in range(n)] for i in range(n)]
        subtours = find_subtours(n, x_val)
        if len(subtours) == 1:
            break
        for st in subtours:
            fs = frozenset(st)
            if fs not in forbidden and len(fs) < n:
                forbidden.append(fs)
    return subtours[0]
\end{lstlisting}

% =============================================================================
\section{Benchmark TSP}

Mesuré sur graphes aléatoires (seed = 42, coordonnées uniformes dans $[0,100]^2$),
solveur OR-Tools CP-SAT, limite de temps 30\,s par instance, CPU.
Les prédictions pour $n \geq 1\,000$ sont extrapolées par loi de puissance
$t = a \cdot n^b$ ajustée sur les mesures réelles ($b \approx 3.18$).

\begin{table}[h]
  \centering
  \resizebox{\textwidth}{!}{%
  \begin{tabular}{cccccc}
    \toprule
    \textbf{Nœuds} & \textbf{\% réussite} & \textbf{\% non-détection}
    & \textbf{\% fausse détection} & \textbf{Temps CPU} & \textbf{Statut} \\
    \midrule
    1        & ---   & --- & --- & $< 0.01$\,s               & INFEASIBLE (trivial)   \\
    10       & 100\% & 0\% & 0\% & $0.02$\,s                 & OPTIMAL                \\
    100      & ---   & --- & --- & $30$\,s (limite)          & FEASIBLE (non optimal) \\
    1\,000   & ---   & --- & --- & $\sim 12$\,h (prédit)     & MTZ hors limites       \\
    10\,000  & ---   & --- & --- & $\sim 2$\,ans (prédit)    & MTZ hors limites       \\
    100\,000 & ---   & --- & --- & $\sim 3000$\,ans (prédit) & MTZ hors limites       \\
    \bottomrule
  \end{tabular}}
  \caption{%
    Résultats benchmark ILP (formulation MTZ).
    \textit{Note n\,=\,100\,:} statut FEASIBLE --- limite de 30\,s atteinte avant preuve
    d'optimalité. \textit{Note n\,$\geq$\,1\,000\,:} la construction du modèle MTZ
    ($n^2$ variables) dépasse les ressources avant même le début de la résolution.
    Utiliser Concorde (DFJ + lazy constraints) pour ces tailles.
  }
\end{table}

\end{document}
